\section{项目的主要内容和技术路线}

\subsection{主要研究内容}
本研究的目标是构建一个面向多源异构知识（Heterogeneous Knowledge）的通用RAG框架，重点解决异构数据融合困难、复杂查询下跨文档证据整合能力不足，以及推理过程缺乏可验证依据等问题。围绕这一目标，本文将研究内容组织为四条主线，并在可复现的端到端原型系统上持续迭代验证其有效性。

在实际应用中，上述问题往往并非独立出现：当查询包含语义歧义、否定约束或需要跨文档拼接线索时，检索结果的噪声会被推理与生成环节放大；而在知识库缺失答案或证据不足的场景下，系统若仍给出高置信回答则会造成更严重的误导。因此，本研究强调从系统视角同时刻画“效果、成本、可信性”三者之间的权衡，并通过可解释的中间结构（证据集合、推理路径、置信度信号）把模型行为落到可分析、可验证的层面。

为保证研究结论具备可对比性，本文将把“检索—排序—融合—生成”作为基础对照管线进行统一评测：检索侧同时覆盖稀疏检索（BM25）与稠密检索（DPR）两类主流范式\cite{bm25, ref6}，排序侧采用交叉编码器进行精排以提升Top-$k$证据质量\cite{bert_rerank}，并在需要多策略集成时引入互惠排序融合（RRF）提升排序稳定性\cite{rrf}。在此基线之上，进一步对比把排序信号显式纳入RAG流程的框架化方法\cite{ref10}与模块化重组思路\cite{ref27}，从而更清晰地区分“单点技巧提升”与“系统机制改进”对最终效果与可信性的贡献。

\textbf{核心问题拆解与递进子任务：}为使研究内容具有可执行性与可验收性，本文将核心问题拆解为四个递进子任务（从“证据入口”到“推理闭环”），并为每个子任务明确阶段目标与输出物。\textbf{子任务1}聚焦多源异构证据统一入口与可追溯定位：面向文本、表格与知识图谱等载体，建立统一证据单元、定位机制与检索算子接口，使检索返回结果能够回溯到原始来源与结构位置，并支持基础实体对齐与来源追溯；高成本的强一致性约束（如版本/时序一致性）作为可选扩展，在工程与数据条件允许时再引入。\textbf{子任务2}面向成本可控的自适应检索与证据选择：在统一入口之上研究“何时检索、检索多少、检索哪些来源”的策略，以证据充分性、答案不一致与约束未满足等信号触发按需检索，并将调用次数、Token消耗与端到端延迟作为硬约束，形成可解释的检索控制与证据筛选机制\cite{ref8, ref12}。\textbf{子任务3}围绕证据链/推理链构造的复杂推理机制展开：在高质量证据集合基础上研究跨文档多跳证据整合与冲突消解，构造可追溯的证据链与推理链；其中可借鉴TRACE\cite{trace}等“基于知识三元组的推理链”思想，把碎片化段落转化为更细粒度的事实单元并进行链式组合，以降低噪声对推理的干扰。\textbf{子任务4}强调可验证推理闭环与系统级评测体系：将推理过程映射为可审计的证据链，设计验证与拒答机制，并建立覆盖效果、成本与证据支撑率的系统级评测闭环，使改进可复现、可归因\cite{ref24, ref25}。上述子任务具有明确递进关系：子任务1提供可覆盖且可定位的证据入口与算子接口；子任务2在入口之上控制证据质量与预算；子任务3在高质量证据上组织复杂推理并输出可追溯链条；子任务4把链条与答案落到可验证与可持续迭代的闭环评测。

\textbf{预期创新点：}
本文预期创新点主要体现在三个方面：其一，面向多源异构载体的证据单元化与可追溯证据空间，提出统一证据单元协议与定位/对齐约束（定位、基础实体对齐、来源追溯），使融合与推理建立在可审计证据之上，而非仅依赖提示拼接；其二，成本约束下的自适应检索控制，把延迟、调用次数与Token预算纳入检索控制信号，并报告效果—成本曲线与预算敏感性，使自适应检索从“有效”走向“可用”\cite{ref12}；其三，证据链驱动的复杂推理与验证闭环，以证据链/推理链作为中间结构组织多跳推理，并将验证与回退机制作为系统控制信号，降低证据不足时的高置信错误输出风险\cite{ref24, ref25}。

\begin{figure}[ht]
\centering
\includegraphics[width=\linewidth]{rag_architecture.pdf}
\caption{检索增强生成（RAG）框架示意图：预训练检索器从外部语料索引中召回Top-$K$证据段落，生成器在输入与证据条件下生成输出，并对不同证据条件下的生成概率进行边缘化处理（改绘自\cite{ref1}）。}
\label{fig:proposal_rag_arch}
\end{figure}

首先，在知识表示与索引层面，面对文本（Text）、知识图谱（KG）、表格（Table）等不同结构的数据形态，单一向量索引难以同时兼顾语义匹配与结构逻辑。考虑到“训练一个同时覆盖文本/图谱/表格的统一嵌入模型”成本高且不稳定，本文采取更可行的折中路线：以“统一证据入口”为目标，强调统一证据单元与检索算子接口，而非强行统一表示空间；在表示层采用结构化序列化与预训练编码器相结合的策略。具体而言，文本侧采用通用语义嵌入并以对比学习增强稳定性\cite{simcse}；表格侧优先采用表格预训练模型（如TAPAS、TaBERT）或将行列与字段约束序列化为自然语言后再进行嵌入，以降低结构差异带来的语义失配\cite{tapas, tabert}；图谱侧以实体描述、关系路径与局部子图摘要作为可检索证据单元，并通过基础实体对齐（如别名归一、ID映射）与路径约束提升证据可追溯性\cite{ref20, ref17}。在系统实现上将多源检索视为“多算子组合”的统一问题，采用“向量+倒排+图遍历”的混合检索算子族，并通过路由与融合策略把它们组织为可复现、可审计的检索过程；其中FAISS\cite{ref32}、Elasticsearch\cite{elasticsearch_guide}与Neo4j\cite{graphdatabases}仅作为算子落地实例，用于分别支持语义近似检索、精确词项匹配与关系遍历，从而为后续复杂推理提供可覆盖、可解释的证据入口。与此同时，本文将结合AutoSchemaKG等方法\cite{ref31}对非结构化文本进行实体关系抽取与Schema归纳，实现面向任务的局部知识图谱构建与可控更新。

在这一主线上，本研究还将关注索引粒度与证据组织方式对系统性能的影响。对于长文档与多主题语料，如何选择分块策略、如何在段落级与文档级之间建立层次化索引，以及如何在多源证据之间维持一致的基础实体对齐与来源标记，都会直接影响检索召回的“可用性”。因此，本文将把索引设计与后续推理需求共同考虑，使检索返回结果不仅“相似”，还尽可能“可组合、可验证”；时间/版本一致性约束作为后续扩展项，在条件允许时再逐步引入。

其次，在检索策略与路由层面，为避免静态检索策略在不同查询类型下产生的效率与噪声问题，本研究将设计面向查询意图的检索控制器（Retrieval Controller）。系统将根据查询的复杂度与意图类型动态选择向量检索、图遍历或混合检索路径，并借鉴Self-RAG\cite{ref8}与主动检索思想\cite{ref12}，在生成或推理过程中以置信度与证据充分性为触发信号实施“按需检索”，使检索过程从“一次性”转向“受控交互式”；对于需要多跳推理的任务，将融合HippoRAG\cite{ref14}与Think-on-Graph\cite{ref17}等思路，在图结构上规划检索路径，并通过重排、并行候选与缓存优先吸收不确定性，只有在关键约束未满足时才触发额外检索，以降低长尾知识定位的搜索成本。

在工程实现上，检索控制器将输出可追踪的决策记录，例如是否触发扩展检索、每一步选择了哪类检索算子以及每步检索的证据贡献，从而支持后续的误差归因与消融分析。通过这种“策略可解释”的设计，本文希望把系统优化从单纯追逐最终指标转向对关键失败模式（歧义误检、约束忽略、证据不足）的针对性修复。

再次，在知识融合与复杂推理层面，检索到的证据往往是碎片化且可能存在冲突的，如何把多源证据组织为可用的推理结构是系统性能的关键。本研究将引入Graph of Thoughts（GoT）\cite{ref5}等图结构推理范式，把推理过程建模为“可分叉、可聚合、可回溯”的思维图，从而更适配多因素制约的复杂问题；在推理闭环上，研究“检索—推理—再检索”的迭代协同机制，但将其严格约束为成本可控的工程形态：设置最大检索轮次与预算上限（如轮次不超过2，且以端到端延迟与Token预算触发提前停止），并优先采用并行候选处理与缓存以避免串行多轮调用导致的不可用耗时。在该约束下，系统仅在关键节点证据不足时触发二次检索，并通过多源证据融合与可信度加权处理冲突信息，提升答案的稳定性与一致性。

考虑到真实证据常常存在冗余、互补与矛盾并存的情况，本研究将把证据融合视为一个可学习或可规则化的“证据选择与对齐”问题，并以基础实体对齐与来源追溯为主来纳入融合过程，避免在开题阶段承诺高成本且难以验收的强对齐与时序版本一致性；当工程与数据条件允许时，再逐步引入更强的一致性约束以提升鲁棒性。在此基础上，推理图结构不仅用于生成中间解释，还用于约束推理过程的证据依赖，从而提升可审计性。

最后，在可解释性与验证层面，为提升系统可用性与可信度，本文将把推理路径映射回知识图谱或证据集合，生成可视化的推理子图，使用户能够直接查看关键实体、关系与证据片段；并引入Chain-of-Verification（CoVe）\cite{ref24}等机制对生成结论进行细粒度自检，结合FactScore\cite{ref25}等指标对事实进行原子级评估，从而把“可解释”从展示层落实到可验证的评估流程。

在具体实现中，本文将把验证机制与无答案/低证据充分性判断结合：当系统检测到证据不足或证据冲突较强时，将优先输出保守结论与证据提示，并在必要时触发二次检索或返回“无法确定”的结果。通过将验证模块前置为系统控制信号，本文希望在保持回答覆盖面的同时，降低高置信错误输出对用户决策的负面影响。

\subsection{技术路线}
本项目将遵循“基础建设—核心突破—系统集成”的研发路径，先完成数据与索引的工程底座，再围绕自适应检索与复杂推理展开算法研发，最后进行系统级集成与评测闭环。第一阶段以数据基础设施与混合索引构建为主：收集Wikipedia、HotpotQA\cite{ref22}与垂直领域PDF、表格等多源数据，完成清洗、去重与分块；利用OpenIE\cite{openie}或基于LLM的抽取工具从文本抽取三元组并构建基础知识图谱；搭建稀疏与稠密索引及图索引组件，并选用成熟的文本/表格/图谱表示方式完成证据单元构建；同时打通“查询—检索—生成”的基础RAG全链路，为后续算法对比提供统一入口。第二阶段面向自适应检索与复杂推理：开发基于分类器或提示的查询路由模块，研究基于置信度信号的主动检索触发策略，并将“调用次数、Token消耗、端到端延迟”作为与效果并列的评测维度，以约束多轮检索推理带来的时间成本；复现并改进GraphRAG与Think-on-Graph等方法，结合节点重要性估计优化图遍历效率；并实现支持分支与合并的GoT推理控制器，形成可维护的思维图状态。第三阶段完成系统集成、优化与评估：将检索、推理与验证模块统一封装为服务或Agent接口，结合缓存与推理加速技术提升响应效率，并在文本多跳数据集（如HotpotQA\cite{ref22}、2WikiMultiHopQA\cite{twiki}）与表格/文本混合数据集（如HybridQA\cite{hybridqa}、OTT-QA\cite{ottqa}、WikiTableQuestions\cite{wikitq}）上开展自动化评测与消融分析，同时在统一多源检索库（文本+表格+局部子图证据单元）上进行\mbox{异构 \& 自适应}端到端评测，辅以人工评估以检验推理路径的合理性与可解释性。

\begin{figure}[ht]
\centering
\setlength{\tabcolsep}{0.8em}
\renewcommand{\arraystretch}{1.2}
\caption{科研与系统实现流程图（本研究）}
\label{fig:proposal_workflow}
\resizebox{0.98\linewidth}{!}{%
\begin{tabular}{c}
\fbox{\begin{minipage}{0.18\linewidth}\centering 文献调研\\问题定义\end{minipage}}
$\rightarrow$
\fbox{\begin{minipage}{0.20\linewidth}\centering 数据准备\\清洗分块/索引\end{minipage}}
$\rightarrow$
\fbox{\begin{minipage}{0.20\linewidth}\centering 方法设计\\检索/重排/推理\end{minipage}}
$\rightarrow$
\fbox{\begin{minipage}{0.20\linewidth}\centering 实验评测\\对比/消融/鲁棒\end{minipage}}
$\rightarrow$
\fbox{\begin{minipage}{0.18\linewidth}\centering 论文撰写\\复现与总结\end{minipage}}\\[0.8em]
\fbox{\begin{minipage}{0.18\linewidth}\centering 系统迭代\\反馈与优化\end{minipage}}
$\leftarrow$
\fbox{\begin{minipage}{0.20\linewidth}\centering 原型实现\\接口/可视化\end{minipage}}
$\leftarrow$
\fbox{\begin{minipage}{0.20\linewidth}\centering 指标记录\\版本管理\end{minipage}}
$\leftarrow$
\fbox{\begin{minipage}{0.20\linewidth}\centering 数据扩展\\难例构造\end{minipage}}
$\leftarrow$
\fbox{\begin{minipage}{0.18\linewidth}\centering 问题回溯\\假设修正\end{minipage}}
\end{tabular}
}
\end{figure}

为确保研究结论具有可复现性与可比性，评测闭环将不仅输出最终答案的正确性指标，还将跟踪检索与推理过程的关键中间量，例如检索召回的覆盖率、重排序的稳定性、推理图结构的规模与分支数量，以及验证模块触发的频率与代价；同时把端到端延迟、LLM调用次数与Token消耗作为与效果并列的约束指标，用于量化自适应检索带来的成本，并在统一推理后端与固定硬件环境下报告p50/p95延迟、不同预算档位下的效果—成本曲线与Pareto最优解。评估设计将采用“分阶段 + 组合验证”的方式以避免结论不可归因：先在纯文本多跳数据集上验证自适应检索/推理机制，再在表格/文本混合数据集上验证异构证据入口与融合能力，最后在统一多源检索库（文本+表格+局部子图证据单元）上引入自适应检索形成\mbox{异构 \& 自适应}系统，并通过消融区分“仅异构”“仅自适应”“异构 \& 自适应”的贡献；其中“异构 + 自适应”指分别验证两类能力的有效性与成本，“异构 \& 自适应”指在同一系统中联合启用异构证据入口与成本约束检索控制。通过把评测指标与系统控制信号联动，本文将能够在“效果提升”和“成本与风险控制”之间做出更清晰的权衡，并据此指导模块级迭代。

\subsection{可行性分析}
从技术基础看，项目组已具备较为完整的工程与算法储备，熟悉Python开发环境，掌握PyTorch\cite{pytorch}、Transformers\cite{huggingface}等深度学习工具链，并能够在现有RAG框架基础上进行模块化开发；同时对向量检索组件（如FAISS\cite{ref32}/Milvus\cite{milvus}）与图数据库（如Neo4j\cite{graphdatabases}）的部署与使用具备实践经验。硬件层面，实验室配备GPU服务器，可支撑重排序模型推理与必要的微调实验。生态层面，Hugging Face提供丰富的开源模型与工具链，GraphRAG、LightRAG等开源项目也为复现与改进提供了可用代码基础。数据层面，HotpotQA、MuSiQue等多跳问答数据集为系统评测提供了可重复、可对比的基准。

更重要的是，本研究已经形成了一个可运行的端到端原型系统：实现了“向量检索—重排序—条件PRF—RRF融合”的完整流水线，并构建了覆盖标准查询、歧义查询、约束查询与无答案查询的测试用例集，能够在统一脚本下输出MRR、NDCG@5与Precision@3等指标\cite{manning_ir, ndcg}并记录迭代轨迹。基于该原型的阶段性实验结果表明，系统在更贴近真实分布的挑战集合上能够暴露问题并快速验证改进方向，这为后续扩大语料规模、引入多模态数据与复杂多跳任务提供了稳定的工程底座与评测闭环。

此外，原型系统已经具备一定的可扩展接口：数据侧能够持续追加文档并重建索引，实验侧能够以迭代编号记录配置与结果，展示侧能够读取历史指标进行可视化对比。这些工程要素使得后续工作可以在“新增数据—新增策略—复现评测”的循环中推进，避免仅停留在概念讨论或一次性实验。基于现有基础，本研究的后续扩展主要集中在两方面：一是将更多异构载体（如表格与图结构）纳入统一索引与证据组织；二是将推理与验证机制从“生成解释”推进到“以证据为约束的系统控制”，从而更贴近高风险应用对可信性的要求。
